% !PRINTED 225
\weberpage{225}
To this it should further be remarked that the superfluously introduced
variables $t_n, t_{n+1} \ldots$ have already been eliminated again in the
second form of $E_{0,s}, E_{1,s} \ldots E_{s-1,s}$, and that the variable
$t_{n-1}$ occurs only in $E_{s,s}$.

\begin{center}
\S.\,73.

\gesperrt{The Bezoutiant.}
\end{center}

Equation (4) of the preceding section, when expanded, will have
the form
\[
y^n + P_1 y^{n-1} + P_2 y^{n-2} + \cdots + P_n = 0,
\]
or, if we determine $t_{n-1}$ so that $S(y) = 0$,
\weq{1}{y^n + P_2 y^{n-2} + \cdots + P_n = 0;}
in it, then,
\[
- 2 P_2 = S(y^2)
\]
is a function of the second degree with respect to the $t$ and with
respect to the $a$. This function is especially important for many
applications and has received from \gesperrt{Sylvester} the name Bezoutiant
of the function $f(x)$, in honour of the French mathematician
\gesperrt{Bezout}, who already in the last century gave the first
correct treatments of elimination.

By the formulas of the last section this function can be
computed comparatively simply. We first introduce, alongside
the variables $t_{n-1}, t_{n-2} \ldots t_0$, a second
system of variables independent of them, $\tau_{n-1}, \tau_{n-2} \ldots \tau_0$, and
set
\weq{2}{
\begin{aligned}
y &= t_{n-1} f_0 + t_{n-2} f_1 + \cdots + t_0 f_{n-1} \\
z &= \tau_{n-1} f_0 + \tau_{n-2} f_1 + \cdots + \tau_0 f_{n-1}.
\end{aligned}
}
Now instead of forming $S(y^2)$, we first compute $S(yz)$,
from which $S(y^2)$ then results when one sets the $\tau$ equal to the $t$.

If one multiplies formula (3) of the preceding section
\weq{3}{y f_s = E_{0,s} f_0 + E_{1,s} f_1 + \cdots + E_{n-1,s} f_{n-1}}
by $\tau_{n-s-1}$ and then sums with respect to $s$ from $0$ to $n-1$,
one finds
\weq{4}{
\begin{aligned}
y z = {} & f_0 \sum_{0,\,n-1}^{s} E_{0,s} \tau_{n-s-1} \\
& + f_1 \sum_{0,\,n-1}^{s} E_{1,s} \tau_{n-s-1} + \cdots + f_{n-1} \sum_{0,\,n-1}^{s} E_{n-1,s} \tau_{n-s-1}.
\end{aligned}
}
